\chapter{Introduction}
\label{ch:intro}

The Large Hadron Collider (LHC) is, as its name suggests, a particle accelerator
designed to collide hadrons at energies far beyond those achieved by previous
experiments. Having recently concluded its third run, it has spent over a decade
probing physics at the fundamental scale. The Standard Model (SM) is currently our
best theory for describing the interactions between elementary particles. It
describes the electromagnetic, weak, and strong interactions, with the latter
taking centre stage throughout this thesis, titled \emph{Jet physics at the LHC
and beyond}. The strong interaction, described by quantum chromodynamics (QCD),
plays a central role in almost every measurement made at the LHC. Whether in
measurements of the Higgs boson, studies of the top quark, or searches for physics
beyond the Standard Model, hard-scattering processes are accompanied by QCD
radiation, giving rise to the complex hadronic final states observed in
experiments. Reliable theoretical predictions for this radiation are therefore
essential for interpreting LHC data and extracting the underlying physics.

Like other quantum field theories, QCD exhibits behaviour that depends on the
energy scale at which it is probed. At high energies, the strong coupling becomes
small, allowing quarks and gluons to interact only weakly. This property, known as
asymptotic freedom, makes perturbation theory applicable and allows systematic
predictions to be obtained as an expansion in the strong coupling, $\alpha_s$. At low energies,
in contrast, the strong coupling grows and perturbation theory eventually breaks
down as quarks and gluons become confined to hadrons. Non-perturbative
methods are then required to describe the physics. Within this thesis, we restrict
ourselves to the high-energy regime where perturbative QCD provides
an excellent description. Nevertheless, obtaining reliable predictions can still present a number of
theoretical challenges.

One such challenge arises from the fact that the quarks and gluons produced in
a hard scattering are not observed directly. Instead, they undergo further QCD
radiation and ultimately hadronise, producing collimated sprays of hadrons known
as \emph{jets}. A jet is an inherently ambiguous object, and defining one
requires a prescription for clustering final-state particles -- whether partons
in a theoretical calculation or energy deposits in a detector -- into jets in a
way that can be applied consistently to both. This prescription is known as a
\emph{jet algorithm}. Different clustering procedures can lead to different jet
configurations, and the choice of algorithm can therefore have a significant
impact on the resulting predictions and on the physical information that can be
extracted. Designing an algorithm suited to a particular process and observable
is consequently a non-trivial task, requiring careful consideration of the
underlying kinematics and the properties one wishes the resulting jets to
possess. This is particularly relevant for Deep Inelastic Scattering (DIS),
where the asymmetric initial state introduces additional features that must be
taken into account.

This motivates the first piece of original work presented in this thesis, described in
Chapter~\ref{chapter:DIS}. Here, we introduce a new inclusive generalised-$k_t$ jet
algorithm for DIS, formulated in the Breit frame, together with an infrared- and collinear-safe
prescription for identifying the jet containing the struck quark. We investigate phenomenological applications of the algorithm, studying the reconstructed jets and assessing their sensitivity to non-perturbative effects such as hadronisation.
The algorithm has been implemented as the \href{https://repo.hepforge.org/source/fastjetsvn/browse/contrib/contribs/DISGenkt}{\texttt{DISGenkt}} plugin to \textsc{FastJet} as part of \texttt{fjcontrib}, and is presented in
Ref.~\cite{vanBeekveld:2026rez}.

A second challenge arises in the form of large logarithms. For many observables,
logarithmically enhanced contributions arise in regions of phase space where the
observable becomes small. These logarithms can become sufficiently large such that
successive terms in the fixed-order perturbative expansion are no longer
parametrically suppressed, reducing the reliability of fixed-order predictions
in these regions. Accurate theoretical predictions therefore require the
logarithmically enhanced contributions to be systematically accounted for at all orders
in perturbation theory through a procedure known as \emph{resummation}.

As with fixed-order calculations, resummation can be performed at successively
higher levels of logarithmic accuracy. Several approaches have been developed to
perform such resummations, and in this thesis we focus on the semi-numerical
\textsc{CAESAR} and \textsc{ARES} formalisms. The second piece of original work
presented in this thesis, described in Chapter~\ref{ch:hhZjet}, extends the
\textsc{ARES} framework to hadronic collisions, incorporating initial-state
radiation at next-to-next-to-leading logarithmic (NNLL) accuracy. The framework
is applied to the process $hh\to Z+\mathrm{jet}$ and completes the crucial conceptual
steps required for its extension to hadronic initial states. The resulting
predictions are validated through a combination of direct and indirect checks,
all of which have been successful.

The remainder of this thesis is structured as follows. Chapter~\ref{ch:QCD}
introduces the theoretical background required for this thesis.
We review the basic structure of QCD, perturbation theory, the running of the
strong coupling, and the description of high-energy scattering processes,
establishing the notation used throughout.

Chapter~\ref{chapter:DIS} presents the generalised-$k_t$ jet algorithm for DIS.
The relevant DIS kinematics and the development of jet algorithms in this
context are introduced before presenting the new algorithm and studying its
phenomenological properties.

Chapter~\ref{ch:resum} provides the background to resummation and
outlines the formalism on which Chapter~\ref{ch:hhZjet} builds.
We discuss the origin of logarithmically enhanced
contributions, the organisation of logarithmic accuracy, and the
\textsc{CAESAR} and \textsc{ARES} approaches to resummation. This provides the
necessary foundation for Chapter~\ref{ch:hhZjet}, where the \textsc{ARES} framework
is extended to hadronic collisions, including the treatment of initial-state
radiation.

Finally, Chapter~\ref{ch:conc} summarises the main results of this thesis and
discusses possible directions for future work and the broader implications of
the research presented here.